\scanpage{153}\mspage{8}%
le \inserted{complexe de} faisceaux $f^!(\Lambda_e)$ (où $f:X\unsure{\text{---}|\longrightarrow}\operatorname{Spec} k=e$) joue le rôle du \uline{complexe des chaines singulières sur} $X$ \uline{à coeff. dans} $\Lambda$, et $Rf_!(f^!{\unsure{\Lambda|\mathbb{\Lambda}}}_e)$ joue le rôle d'une \uline{homologie} $H_*(X/e)$, vis à vis bien sûr de coefficients sur $e$ qui soient des complexes de $\Lambda$-modules. Tu justifieras aisément cette assertion à l'aide des ``th. de dualité globale'', par un ou deux tours de passe-passe que je t'épargne ici.

\uline{Remarques} 1) Il n'y a pas à tronquer, ça marche en toutes dimensions

2) C'est (au moins du p\textsuperscript{t} de vue philosophique) lié au fait qu'on dispose pour les coefficients envisagés \inserted{(sous condition de ``constructibilité\unsure{'')|)}} d'un th. de ``bidualité''\deleted{\unsure{)}}, du moins sous réserve de résolution des singularités (mais Deligne m'a dit je crois qu'il sait le prouver sans ça) à valeurs dans un ``complexe dualisant'' $K_e$ \inserted{(sur $e$)} $\unsure{\text{--}|;|\text{--};}$ $K_X$ (sur $X$). Si p.ex. $\Lambda$ est ``\unsure{autodual}'' \inserted{(ou Gorenstein)} p.ex. $\Lambda=\mathbb{Z}/n\mathbb{Z}$, on peut prendre $K_e=\Lambda$, donc \inserted{le complexe dualisant} $K_X=f^!(K_e)$ n'est autre que le ``complexe des chaines singulières à \unsure{coeff.} \unsure{dans} $\unsure{\Lambda}$'' \illegible{1}
% archive sheet number: 153
